\documentclass[12pt]{article}
\usepackage[utf8]{inputenc}
\usepackage[T1]{fontenc}
\usepackage[margin=1in]{geometry}
\usepackage{setspace}
\usepackage{xcolor}
\usepackage[colorlinks=true, urlcolor=blue!60!black, linkcolor=blue!60!black]{hyperref}
\setlength{\parindent}{2em}
\setlength{\parskip}{1.25em}
\renewcommand{\baselinestretch}{1.0}
\begin{document}
\begin{center}
{\LARGE \textit{United States v.\ Skrmetti}, 605 U.S.\ 495 (2025)\par}
\vspace{0.8em}
{\large Case Brief --- Four-Page Edition\par}
\end{center}
\vspace{1em}

\section{Memory Jogger}

Tennessee's ban on puberty blockers and hormone therapy for minors with gender dysphoria classifies by age and medical use, not sex or transgender status, so it gets only rational basis review and is constitutional (\href{https://www.courtlistener.com/opinion/10776831/united-states-v-skrmetti/}{CourtListener}).

\section{Detailed Case Facts}

In March 2023 Tennessee enacted SB1, barring providers from giving minors puberty blockers or hormones to enable them to live as an identity inconsistent with their biological sex or to treat distress from that discordance. The same drugs remain available to minors for a congenital defect, precocious puberty, disease, or injury, but ``gender dysphoria, gender identity disorder, [and] gender incongruence'' are excluded from those categories. The legislature found risks of irreversible sterility, disease, and regret, noted restrictions in Sweden, Finland, and the United Kingdom, and declared an interest in ``encouraging minors to appreciate their sex.'' SB1 is enforced by the Attorney General, licensing discipline, and a private right of action. Three transgender adolescents, their parents, and a physician sued; the United States intervened, and after the change in administrations it told the Court in February 2025 that SB1 is constitutional.

\section{Procedural History}

The district court found no standing to challenge the surgery ban but preliminarily enjoined the medication ban, holding that transgender persons are a quasi-suspect class and that SB1 likely failed intermediate scrutiny. The Sixth Circuit (Chief Judge Sutton; Judge White dissenting) reversed, applying rational basis review. The Supreme Court granted the United States' petition limited to equal protection and affirmed on June 18, 2025.

\section{Judicial Votes}

Affirmed, \textbf{6--3}.
\begin{itemize}
\item \textbf{Majority:} Chief Justice Roberts, joined by Justices Thomas, Gorsuch, Kavanaugh, and Barrett, and by Justice Alito as to Parts I and II-B.
\item \textbf{Concurrences:} Justice Thomas; Justice Barrett, joined by Thomas; Justice Alito (in part and in the judgment).
\item \textbf{Dissents:} Justice Sotomayor, joined by Jackson and (except Part V) Kagan; Justice Kagan.
\end{itemize}

\section{Holding}

A law barring puberty blockers and hormones for minors' gender dysphoria while permitting them for other diagnoses classifies by age and medical use, neither of which triggers heightened scrutiny, and it is rationally related to protecting minors amid medical uncertainty. Even assuming \textit{Bostock} applies beyond Title VII, sex is not a but-for cause of SB1's operation. The Court did not decide whether transgender status is a suspect class, did not reach the surgery ban or parental rights, and noted (footnote 3) that a law targeting a class of persons, rather than treatments, would differ.

\section{Analysis of Opinions}

\textbf{Majority (Roberts).} Age and medical-use lines get rational basis review (\textit{Murgia}, \textit{Vacco v.\ Quill}). A medical treatment is both drug and indication, so puberty blockers for gender dysphoria differ from those for precocious puberty; unlike \textit{Loving}, SB1 bars nothing for one sex that it allows the other, and the ``appreciate their sex'' finding shows no invidious purpose under \textit{Feeney}. Under \textit{Geduldig}, SB1 does not classify by transgender status, because the excluded group includes only transgender minors while the permitted group includes both. Changing a minor's sex would not supply a qualifying diagnosis, so \textit{Bostock}'s logic fails. Citing sterility risks and the 2024 Cass Review, the Court applied \textit{Gonzales v.\ Carhart}'s latitude amid medical uncertainty.

\textbf{Thomas (concurring).} \textit{Bostock} should never enter equal protection, and courts should not defer to ``self-described experts,'' citing WPATH's response to political pressure.

\textbf{Barrett (concurring).} Transgender status is not a suspect class: it is not an ``obvious, immutable, or distinguishing characteristic,'' the group is ``large, diverse, and amorphous,'' and future inquiries should require a history of \textit{de jure} discrimination.

\textbf{Alito (concurring in part).} There is a ``strong argument'' that SB1 classifies by transgender status, but that class is not suspect.

\textbf{Sotomayor (dissenting).} SB1 classifies by sex on its face: a girl with excess facial hair may receive testosterone suppressors, but a boy with the same distress may not. Transgender people are at least quasi-suspect, and SB1 would fail intermediate scrutiny. \textbf{Kagan (dissenting)} would have remanded for heightened scrutiny without predicting the result.

\section{Examples: Future Applications}

\begin{enumerate}
\item \textbf{Same-Side: A Parallel Ban.} Another state's ban on puberty blockers for gender dysphoria with licensing penalties gets rational basis review.
\item \textbf{Same-Side: Medicaid Exclusion.} Excluding coverage of hormones and surgery for gender dysphoria, while covering them for cancer, would likely be a diagnosis-based line under \textit{Geduldig}.
\item \textbf{Opposite-Side: Dress Codes.} A rule barring clothing ``inconsistent with'' biological sex classifies by sex and triggers intermediate scrutiny.
\item \textbf{Opposite-Side: Status-Based Insurance Denial.} Denying all coverage to anyone who ``identifies as transgender'' targets persons, outside the footnote 3 distinction.
\item \textbf{Fence-Sitter: Birth Certificates.} Barring minors from changing sex markers may be framed as recording biology or as classifying persons; the Court vacated and remanded \textit{Stitt v.\ Fowler} in light of \textit{Skrmetti}.
\end{enumerate}

\section{Critique}

\textbf{Conservative and originalist.} Sarah Parshall Perry praised the Court for leaving pediatric medicine to voters and refusing to import \textit{Bostock}; scholars see Barrett's \textit{de jure} test as an objective anchor. Yet Alito questioned the \textit{Geduldig} move, and some find the case hard to reconcile with \textit{Bostock}.

\textbf{Progressive.} Kate Redburn argues the Court created ``double deference'' in an undefined ``medical context,'' and critics say reviving \textit{Geduldig} lets legislatures shield sex classifications by drafting around diagnoses.

\section{Key Quotations}

\begin{quote}
``In the medical context, the mere use of sex-based language does not sweep a statute within the reach of heightened scrutiny.''
\end{quote}
\noindent --- Chief Justice Roberts, for the Court.

\begin{quote}
``The Equal Protection Clause does not resolve these disagreements. Nor does it afford us license to decide them as we see best. Our role is not `to judge the wisdom, fairness, or logic' of the law before us, \textit{Beach Communications}, 508 U.\ S., at 313, but only to ensure that it does not violate the equal protection guarantee of the Fourteenth Amendment. Having concluded it does not, we leave questions regarding its policy to the people, their elected representatives, and the democratic process.''
\end{quote}
\noindent --- Chief Justice Roberts, for the Court.

\begin{quote}
``The Court's willingness to do so here does irrevocable damage to the Equal Protection Clause and invites legislatures to engage in discrimination by hiding blatant sex classifications in plain sight. It also authorizes, without second thought, untold harm to transgender children and the parents and families who love them.''
\end{quote}
\noindent --- Justice Sotomayor, dissenting.

\end{document}
